% status: SKELETON
% Chapter 39 of the second edition. Plan: docs/STRUCTURE.md. Rules: AUTHORING.md.
\chapter{Benchmarking Without Lying to Yourself}\label{ch:benchmarking}

\begin{ChapterIntent}
A benchmark is an experiment, and most inference benchmarks are poorly
controlled ones. This chapter defines the metrics that matter --- time to
first token, inter-token latency, throughput and goodput --- and how to
measure them: open- and closed-loop load, realistic arrival and length
distributions, warm-up and cache state, percentiles rather than means, and
comparisons that hold model, quantization and stopping rules fixed. It ends
with a record format that makes a number reproducible, and the common ways
results mislead.
\end{ChapterIntent}

\OpeningSection{The same engine, three times apart}

\NapkinSection{Napkin math: how many runs are enough}

\section{What to measure}

\section{Open and closed loops}

\section{Arrivals and workloads}

\section{Warm-up, caches and steady state}

\section{Percentiles and variance}

\section{Comparing engines fairly}

\section{Industry benchmarks}

\section{A benchmark record}

\LedgerSection

\LabSection{one server, two load generators}

\HappenedSection{measurements that compared a system with itself}

\FrontierSection{}

\ExercisesSection
